% !PRINTED 605
\weberpage{605}
Umgekehrt bilden in diesem Falle 1) die $p$ rationalen Zahlen $0, 1, 2, \ldots, p-1$ ein volles Restsystem nach dem Modul $\pi$, und es ist folglich jede Zahl in $\mathfrak{o}$ und mithin auch $\Theta$ mit einer rationalen Zahl congruent. Also haben wir

1) \quad $N_a(\pi) = p$ \quad $a_{2,2} = 1$

2) \quad $N_a(\pi) = p^2$ \quad $a_{2,2} = p.$

Im ersten Falle ist $a_{1,2}$ aus der Bedingung

\weq{3}{a_{1,2} + \Theta \equiv 0 \pmod{\pi}}
zu bestimmen, im zweiten kann $a_{1,2} = 0$ genommen werden. Die Basisform von $\pi$ ergiebt sich also, wenn $u, v$ die Variablen sind,

1) \quad $\pi = pu + (a_{1,2}+\Theta)v$

2) \quad $\pi = p(u+\Theta v),$

woraus man sieht, dass $\pi$ im Falle 2), wie zu erwarten war, ein \gesperrt{Hauptfunctional} ist.

Um im Falle 1) die Zahl $a_{1,2}$ zu bestimmen, bilden wir die Norm von $a_{1,2}+\Theta$, die durch $p$ theilbar sein muss. Wenn $\Theta$ durch Aenderung des Vorzeichens von $\sqrt{d}$ in $\Theta'$ \"ubergeht, so ist nach \S.\,164, (8) $\Theta - \Theta' = \sqrt{\Delta}$, und f\"ur die Norm von $a_{1,2}+\Theta$ ergiebt sich nach \S.\,164 (6), a), b):

\weq{4}{N(a_{1,2}+\Theta) = (a_{1,2}+\Theta)(a_{1,2}+\Theta')}
\[ {} = a_{1,2}^2 - d, \quad d\equiv 2, 3\ (\mathrm{mod}\ 4) \]
\[ {} = a_{1,2}^2 + a_{1,2} + \frac{1-d}{4}, \quad d\equiv 1\ (\mathrm{mod}\ 4). \]

Wenn diese Zahl durch $p$ theilbar ist, so ist einer der beiden Factoren $a_{1,2}+\Theta$, $a_{1,2}+\Theta'$ durch $\pi$ theilbar. Da aber $\Theta' = -\Theta$ oder $=1-\Theta$ ist, also

a) \quad $a_{1,2} + \Theta' = -(-a_{1,2}+\Theta)$

b) \quad $a_{1,2} + \Theta' = -(-a_{1,2}-1+\Theta)$

ist, so kann man $a_{1,2}$ immer so annehmen, dass gerade $a_{1,2}+\Theta$ durch $\pi$ theilbar wird, und $\Theta$ ist also wirklich mit einer rationalen Zahl congruent. Die nothwendige und hinreichende Bedingung daf\"ur, dass $p$ zwei Primfactoren enth\"alt, ist folglich die, dass $a_{1,2}^2 - d$ oder $a_{1,2}^2 + a_{1,2} + \frac14(1-d)$ durch geeignete Annahme \"uber $a_{1,2}$ durch $p$ theilbar wird. Durch Multiplication mit 4 gehen diese beiden Bedingungen \"uber in
\[ 4a_{1,2}^2 - 4d \equiv 0, \quad (2a_{1,2}+1)^2 - d \equiv 0 \pmod{4p}, \]
%CONT>>
